\section{People's Hall}
\label{sec:peoplesHall}

If the characters attempt to meet \npcTroublemaker{} the next morning, start the scene with \fullref{ssec:peoplesHallCrowd}.

If they decide to skip this meeting, finish this chapter and proceed to \nameref{chap:chapter2}. There, they can learn what happened with \npcRebelArdent{} while hearing the gossiping people.

\subsection{Gathering Crowd}
\label{ssec:peoplesHallCrowd}

Start the scene by reading the following:

\quoteBox{
	As you approach the market and the People's Hall, you can see a small crowd rushing inside the Hall. The food merchant you wanted to trade with just closed his store. "Is that true? I must see it!", he mutters to himself as he rushes next to you.
}

As the characters follow the crowd inside, proceed to \fullref{ssec:peoplesHallGlyphs}.




\subsection{Ten Glyphs}
\label{ssec:peoplesHallGlyphs}

Once the PCs push through the crowd, read the following:

\quoteBox{
	You stop a few paces before a dais with the Common Throne, beneath which there is an ardent kneeling. Around you can see ten glyphs, one for each foolish attribute: \\ \\
	Avarice \\
	Callousness \\
	Profligacy \\ \\
	and so on...
}

A PC with an expertise in Vorin theology will recognize the glyphs as Ten Foolish Attributes from the teachings.

If at least one of the characters can read alethi women's script, they will be able to read:

\quoteBox{
	Below each of the glyphs there is a smaller writing: \\ \\
	Avarice \\
	The queen gathers goods she does not have a use for. \\ \\
	Callousness \\
	Blind to her people's needs, the queen terminated Beggar's Feast. \\ \\
	Profligacy \\
	Food rots and is wasted in the palace, while the city starves.	
}

\spoilerBox{
	Events leading to this moment are described in \bookWOR{}, Interlude I-12: Lhan.
}






\subsection{Blind Accusations}

After the PCs discover the glyphs, read the following:

\quoteBox{
	Suddenly several guards with queen's colors enter the People's Hall. One of them approaches you. \\ \\
	"Who did this?", he asks. "Was it you? Hey answer me! Why did you do that!", he shouts, his arm reaching to his sword.
}

\enemyQueensGuard{} accuses one of the PCs of painting the glyphs. With hand on his sword, he is ready to arrest the characters, drawing weapons on the first sign of resistance. More \enemyQueensGuard{}s are nearby - starting a fight will probably make them join it.

The characters may try to talk their way out of the trouble. Run it as a conversation.
\conversationDefence{\enemyQueensGuard{}}{15}{18}{5}
The guard is clearly agitated and blind to reason. He fears the queen will punish him unless he quickly finds a person responsible for the glyphs.

In the conversation, \enemyQueensGuard{} will use mostly \skillIntimiation{}.
\conversationEnd{\enemyQueensGuard{}}{1}

\reasonBox{
	Try to convince your players the situation is serious. Make them terrified of a huge guard in full armor and his companions. Make other NPCs scared of the guards, making it clear they will not join the fight.
}

\reasonBox{
	In this scene, there is no danger for the PCs. Instead, use it to make them roleplay their characters in the dire situation. What will they do? Will they instantly accuse harmless ardent? Or will they make a heroic sacrifice to protect her?
}

\subsection{\npcRebelArdent{}'s Capture}

Regardless of the conversation result, follow with the below:

\quoteBox{
	"Hey, give it a rest!", another guard says. "It's clearly the ardent, those here know nothing about it. You two, take her away", he adds, pointing at the kneeling clerk. \\ \\
	When the ardent is taken away, there is no fear in her eyes. Only calm acceptance. 
}

Once the guards take \npcRebelArdent{} away, you can complete this chapter.